%!TEX program = lualatex
\documentclass[11pt]{article}

% ---------- Font setup ----------
\usepackage{fontspec}
\usepackage{unicode-math}
\setmainfont{Libertinus Sans}[Scale=1.1]
\setmathfont{Libertinus Math}[Scale=1.1]

% ---------- Page layout ----------
\usepackage[a4paper,margin=1in]{geometry}

% ---------- Line spacing ----------
\usepackage{setspace}
\setstretch{0.96}
\usepackage[colorlinks=true, urlcolor=blue]{hyperref}

% ---------- Packages ----------
\usepackage{tikz-cd}
\tikzset{
  no line/.style={draw=none,
    commutative diagrams/every label/.append style={/tikz/auto=false}}}

% ---------- Table of contents -------
\hypersetup{linkcolor=black}
\renewcommand{\partname}{Week}

% ---------- Environments ---------
\usepackage{amsthm}
\theoremstyle{definition}
\newtheorem{theorem}{Theorem}[part]
\newtheorem{proposition}[theorem]{Proposition}
\newtheorem{lemma}[theorem]{Lemma}
\newtheorem{corollary}[theorem]{Corollary}
\newtheorem{definition}[theorem]{Definition}
\newtheorem{example}[theorem]{Example}
\theoremstyle{remark}
\newtheorem{remark}[theorem]{Remark}
\renewcommand{\thetheorem}{\arabic{theorem}}

% ---------- Custom commands ----------
\newcommand{\Lan}{\operatorname{Lan}}
\newcommand{\Ran}{\operatorname{Ran}}
\newcommand{\Ho}{\operatorname{Ho}}
\newcommand{\Hom}{\operatorname{Hom}}
\newcommand{\id}{\operatorname{id}}
\newcommand{\im}{\operatorname{im}}
\newcommand{\colim}{\operatorname{colim}}
\newcommand{\Top}{\mathbf{Top}}
\newcommand{\Ch}{\mathbf{Ch}}
\newcommand{\Ab}{\mathbf{Ab}}
\newcommand{\ChAb}{\mathbf{Ch}(\mathbf{Ab})}
\newcommand{\bZ}{\mathbb Z}
\newcommand{\calC}{\mathcal{C}}
\newcommand{\calD}{\mathcal{D}}
\newcommand{\calE}{\mathcal{E}}
\newcommand{\calF}{\mathcal{F}}
\newcommand{\calM}{\mathcal{M}}
\newcommand{\calN}{\mathcal{N}}
\newcommand{\calV}{\mathcal{V}}
\newcommand{\calW}{\mathcal{W}}
\newcommand{\bfL}{\mathbf{L}}
\newcommand{\bbL}{\mathbb{L}}
\newcommand{\bfset}{\mathbf{Set}}

% ---------- Title ----------
\title{
    Reading Group on Categorical Homotopy Theory
}
\author{}
\date{\today}

\begin{document}

\maketitle
\tableofcontents
\newpage

\part{Motivation --- Mar 3, led by Aref}

\section{To invert weak equivalences}
In the (convenient) category $\Top$ of topological spaces,
homotopy equivalence is a good equivalence relation,
so we can define the homotopy category $\Ho\Top$,
whose objects are all the topological spaces
but morphisms are equivalence classes of continuous maps up to homotopy equivalence.
Sometimes working in $\Ho\Top$ is easy since every homotopy equivalence becomes an isomorphism. 

There are also other interesting notions of ``weak equivalences'' that are not actually equivalences.
Again, in the category $\Top$ of topological spaces, we say a map $f : X \to Y$ is a
\href{https://ncatlab.org/nlab/show/weak+homotopy+equivalence}{weak homotopy equivalence} if
$$\pi_n(f) : \pi_n(X) \to \pi_n(Y)$$
is an isomorphism for $n\ge 1$ and a bijection when $n = 0$.
Clearly homotopy equivalence is a weak equivalence, and we have the classical Whitehead's theorem
(which is not provable in HoTT!) saying that,
every weak homotopy equivalence between CW-complexes is a homotopy equivalence.
However, weak homotopy equivalences are generally not invertible
(see \href{https://ncatlab.org/nlab/show/pseudocircle}{here} for a counterexample).

Another important example comes from the category of
\href{https://ncatlab.org/nlab/show/chain+complex}{chain complexes}.
\href{https://ncatlab.org/nlab/show/quasi-isomorphism}{Quasi-isomor-phisms},
as the name might indicate, are also generally not invertible.

What we can do is formally add an invertible arrow to theses weak equivalences.
We consider the following construction of $\Ho\Top_{whe}$.
The objects are all the topological spaces.
The morphisms between $X$ and $Y$ are the equivalence classes of finite zig-zags
$$X \leftarrow Z_1 \rightarrow \cdots \leftarrow Z_{n} \to Y,$$
suject to
\begin{enumerate}
    \item all the arrows backwards are weak homotopy equivalences.
    \item identity arrows can be removed and then composed.
\end{enumerate}

Then a weak homotopy equivalence $X \stackrel{w}{\to} Y$ has an inverse given by
$$Y \stackrel{w}{\leftarrow} X \stackrel{\id}{\to} X.$$

However, this construction may not give a locally small category.
That is, given a category $\mathcal C$ and a class of morphisms $W$ we want to want to invert,
the resulting $\Ho_W \mathcal C(X,Y)$ may be a proper class
(or in type theory, this means the type would lie in the next universe level).

This is one of the reasons why we need
\href{https://ncatlab.org/nlab/show/model+category}{model category}.
With auxiliary data of fibrations and cofibrations,
we can ensure the homotopy category we obtain is indeed locally small.

This process is also called \href{https://ncatlab.org/nlab/show/localization}{localization},
which rougly mean to ignore weak equivalences and focus ``locally'' on other maps.
We will call the canonical functor $\gamma : \calC \to \Ho \calC$ the localization functor.

\section{Extending functors over localization}
Given two categories $\calC$ and $\calD$ and a functor $F : \calC \to \calD$,
we may want to extend $F$ to $\Ho\cal C \to \Ho\cal D$.
That is, we want $\tilde F$ such that the following diagram commutes
\[\begin{tikzcd}
\calC \arrow[r,"F"] \arrow[d,"\gamma"'] & \calD \arrow[d,"\gamma"] \\
\Ho \calC \arrow[r,"\tilde{F}"'] & \Ho \calD
\end{tikzcd}\]

This is easy if $F$ preserves weak equivalences.
But unfortunately, this is most of the time not the case.
For example, the following two diagrams are (component-wise) homotopy equivalent

\[
\begin{tikzcd}
    D ^ 1 \arrow[d] & S ^ 0 \arrow[l] \arrow[r] \arrow[d] & D ^1 \arrow[d] \\
    * & S ^ 0 \arrow[l] \arrow[r] & *
\end{tikzcd}
\]
But the colimit (a functor $\Top^{\bullet \leftarrow \bullet \rightarrow} \to \Top$) of the first row
is $S^1$, while the colimit of the second row is $*$. 

So we have to compromise by considering a weaker notion of extension,
i.e., Kan extension of the functor $F$ along $\gamma: \calC \to \Ho\calC$.
This gives us the so-called derived functors.
In particular, homotopy co/limits are then the derived functors of (strict) co/limits.

\section{``Weak equivalences are all that matter''}
Traditionally, all the constructions mentioned above take place in model categories.
But a generalization is possible: we can introduce a framework
for constructing derived functors between categories equipped with a reasonable notion of
weak equivalence. As Riehl said in the preface:

\begin{quote}
``The weak equivalences are all that matter in abstract homotopy theory, showing
that particular notions of cofibrations/fibrations and cofibrant/fibrant objects are irrelevant
to the construction of derived functors.''
\end{quote}

This generalization allows us to produce derived functors where
appropriate model structures for the diagrams do not necessarily exist.

Also, this philosophy could be inspiring for type theorists.
In type theories, we rarely have a nice model structure,
but we always have a nice notion of equivalence.
So we would at least hope studying this book can help us understand
homotopy co/limits in type theories like HoTT,
or compare homotopy co/limits and strict co/limits in 2LTT.
%In 2.2, the book introduces this notion:
%a left deformation on $\calC$ is an endofunctor $Q$
%with a weak equivalence $Q \Rightarrow 1$.
%Doesn't it resemble modality?
%We hope they can be somehow connected.

\section{A generalization of Whitehead theorem}

Now that we have introduced the notion of localization,
we can reinterpret the following statements in classical algebraic topology:

Every homotopy equivalence is a weak homotopy equivalence $\longrightsquigarrow$
The localization of weak homotopy equivalence factors through the localization of homotopy equivalence, i.e,

\[
\begin{tikzcd}
    & \Top \arrow[dl] \arrow[dr] & \\
    \Ho \Top_{he} \arrow[rr, dashed] &  & \Ho \Top_{whe}
\end{tikzcd}
\]

Whitehead theorem, i.e., every weak homotopy equivalence between CW complexes is a homotopy equivalence
$\longrightsquigarrow$
The dashed functor above restricts to an equivalence on the full subcategory of CW complexes.

It turns out that $\Ho \Top_{he}$ also arises directly from the enrichement setting.
Let $\underline{\Top}$ denote $\Top$ enriched over itself.
Define $h\Top$ to be the category with objects the same as $\Top$ but
$$h\Top(X,Y) := \pi_0\underline{\Top}(X,Y).$$
Note that $h\Top$ is equivalent to $\Ho\Top_{he}$, and hence
the statements above can be seen as sort of a comparison between
two notions of ``homotopy category'' arising from different settings: the enriched one and the homotopical one.

One of the main goals of this reading group is to understand Theorem 10.5.1,
which generalizes the observation mentioned above:

Let $\calV$ be a category with some structures and let $\underline{\calM}$
be a $\calV$-homotopical category (whatever this means). Then there is a
$\Ho\calV$-functor $\varphi : \underline{h\calM} \to \underline{\Ho\calM}$
(however $h\calM$ is defined) whose underlying unenriched functor commutes
with both localizations
\[
\begin{tikzcd}
    \calM \arrow[d] \arrow[dr] & \\
    h\calM \arrow[r, "\phi"] & \Ho\calM
\end{tikzcd}
\]
and the restriction of $\varphi$ to fibrant-cofibrant objects (whatever this means)
is fully faithful.

\newpage
\part{Derived functors, featuring Kan extensions and homological algebra --- Mar 10, led by Mark}
\setcounter{section}{0}
\section{Homotopical category}
\begin{definition}
    (Def 2.1.1)
    A \emph{homotopical category} is a category $\calM$ with a wide subcategory $\calW$ (containing all objects of $\calM$)
    satisfying 2-of-6 property:
    for any composable triple of arrows $f,g,h$,
    if $hg,gf\in \calW$, then $f,g,h,hgf \in \calW$.
    The arrows in $\calW$ are called weak equivalences.
\end{definition}

\begin{remark}
    2-of-6 implies 2-of-3.
\end{remark}

\begin{definition}[Def 2.1.6]
    The \emph{homotopy category} $\Ho\calM$ of a homotopical category $(\calM, \calW)$
    is the formal localization of $\calM$ at the subcategory $\calW$. 
\end{definition}

\begin{remark}
    The localization is obtained by formally inverting morphisms in $\calW$.
    It is not guaranteed to be locally small, but we put this issue aside for now.
\end{remark}

\begin{proposition}
    \label{universal-property-localization}
    There is a localization functor $\gamma : \calM \to \Ho\calM$ with
    the universal property that for any $F : \calM \to \calN$ sending weak equivalences to isomorphisms
    factors uniquely through $\gamma$:

    \[
    \begin{tikzcd}
        \calM \arrow[d, "\gamma"] \arrow[dr, "F"] & \\
        \Ho\calM \arrow[r, dashed] & \calN
    \end{tikzcd} 
    \]

\end{proposition}

\begin{definition}
    A functor between two homotopical categories is \emph{homotopical} if it preserves weak equivalences.
\end{definition}

\begin{remark}
    If $F : \calM \to \calN$ is homotopical, then by the Proposition \ref{universal-property-localization},
    there exists a unique functor $F : \Ho\calM \to \Ho\calN$,
    such that the following square commutes
    \[
    \begin{tikzcd}
        \calM \arrow[d, "\gamma"] \arrow[r, "F"] & \calN \arrow[d, "\gamma"]\\
        \Ho\calM \arrow[r, dashed, "F"] & \Ho\calN
    \end{tikzcd} 
    \]
\end{remark}

Singular (co)homology chain complex functor from $\Top$ to $\ChAb$
(see Section \ref{Homological-algebra}) is homotopical.
It's nice, but in practice, functors of interests are rarely homotopical.
Counterexamples are $\colim$, e.g., (strict) pushout of spaces, and $-\otimes A : \ChAb \to \ChAb $
for most of abelian groups $A$.
For this reason, we introduce homotopical approximations.
Instead of requiring a strict commutativity, we allow lax transformations
comparing the original functor and its approximation.

\begin{definition}(Def 2.1.17)
    Let $F : \calM \to \calN$ be a functor between homotopical categories.
    A \emph{total left derived functor} $\bfL F$ is a right Kan extension
    (see Section \ref{Kan-extension}) of $\delta F$ along
    $\gamma$:
    \[
    \begin{tikzcd}
    \calM \arrow[d,"\gamma"] \arrow[rr,"F"] 
    & {} 
    & \calN \arrow[d,"\delta"] \\
    \Ho\calM \arrow[rr,dashed,"\bfL F"'] 
    & {} \arrow[u,Rightarrow,"\alpha"',shorten <=6pt,shorten >=6pt]
    & \Ho\calN
    \end{tikzcd}
    \]
\end{definition}

Note that $\bfL F$ is a functor between homotopy categories,
In practice, it is better not to work only on homotopy categories,
since they can be badly behaved—for instance, they often have very few (ordinary) colimits.
Instead, we would like to consider the lift of $\bfL F$ to the original categories.

\begin{definition}(Def 2.1.19)
    A \emph{left derived functor} of $F : \calM \to \calN$ is a homotopical functor
    $\bbL F : \calM \to \calN$ with a natural transformations $\lambda : \bbL F \to F$
    such that $(\delta\bbL F, \delta\lambda)$ is a total left derived functor of $F$.
    \[
    \begin{tikzcd}
        \calM \arrow[d,"\gamma"]
            \arrow[rr,"F", bend left=20]{r}[name=UU]{}
            \arrow[rr,"\bbL F"', bend right=20]{r}[name=UD]{}
        &
        \arrow[Leftarrow, to path = (UU) -- (UD)\tikztonodes]{}{\lambda}
        & \calN \arrow[d,"\delta"] \\
        \Ho\calM \arrow[rr,dashed,"\sigma\bbL F"'] 
        & {}
        & \Ho\calN
    \end{tikzcd}
    \]
\end{definition}

\section{Kan extension}\label{Kan-extension}
Intuitively, Kan extension of $F : \calC \to \calE$ along $K : \calC \to \calD$
is a best approximation of the lifting problem
\[
\begin{tikzcd}
    \calC \arrow[d, "K"] \arrow[r, "F"] & \calE \\
    \calD \arrow[ur, dashed, "?"'] &  
\end{tikzcd}
\]
which usually does not have a strict solution.

\begin{definition}(Def 1.1.1)
    Given functors $F : \mathcal{C} \to \mathcal{E}$, $K : \mathcal{C} \to \mathcal{D}$, a \emph{left Kan extension} of $F$ along $K$ is a functor 
    \[\Lan_K F : \mathcal{D} \to \mathcal{E}\] together with a natural transformation
    \[\eta : F \Rightarrow (\Lan_K F)\circ K\]
    such that for any other such pair $(G : \mathcal{D} \to \mathcal{E},\; \gamma : F \Rightarrow GK)$,
    $\gamma$ factors uniquely through $\eta$.

    Dually, a \emph{right Kan extension} of $F$ along $K$ is a functor
    \[\mathrm{Ran}_K F : \mathcal{D} \to \mathcal{E}\]
    together with a natural transformation
    \[\epsilon : (\mathrm{Ran}_K F)\circ K \Rightarrow F\]
    such that for any $(G : \mathcal{D} \to \mathcal{E},\; \delta : GK \Rightarrow F)$,
    $\delta$ factors uniquely through $\epsilon$.
\end{definition}

By the Yoneda lemma, we see the the universal property of left Kan extensions
is equivalent to the assertion that
$(\Lan_K F, \eta)$ represents the functor $\calE ^{\calC}(F, -\circ K) : \calE ^{\calD} \to \bfset$.
Dually, $(\Ran_K F, \epsilon)$ corepresents the functor $\calE ^{\calC}(-\circ K, F) : \calE ^{\calD} \to \bfset$.

We further observe that if the left and right Kan extensions
of all functors from $\calC$ to $\calE$
along a fixed functor $K : \calC \to \calD$ exist,
then $\Lan_K$ and $\Ran_K$ encode the left and right adjoints
of $K^* : \calE ^{\calD} \to \calE ^{\calC}$ respectively.

\[
\begin{tikzcd}
    {\calE}^{\calC} \arrow[rr, "\Lan_K"{name = Lan, description}, bend left=45]
    \arrow[rr, "\Ran_K"{name = Ran, description}, bend right=45]
    & & {\calE}^{\calD} \arrow[ll, "K^*"{name = K, description}]
    \arrow[phantom, from=K, to=Lan, "\dashv" rotate=-90, no line]
    \arrow[phantom, from=Ran, to=K, "\dashv" rotate=-90, no line]
\end{tikzcd}
\]

There are so many example of extensions. See ``1.4 All concepts'' for some of them justifying the
slogan that ``All concepts are Kan extensions''.

In some good settings, Kan extensions always exist and can be computed concretly:
\begin{theorem}(Th 1.2.1)
    When $\calC$ is small, $\calD$ is locally small, and $\calE$ is cocomplete,
    the left Kan extension of any functor $F : \calC \to \calE$ along any functor
    $K : \calC \to \calD$ is computed at $d \in \calD$ by the colimit
    \[\Lan_K F(d)= \int^{c \in \mathcal{C}} \coprod_{i \in \mathcal{D}(Kc,d)} Fc\]
    and in particular necessarily exists.
\end{theorem}

Here $\int^{\calC}H$ is the coend for a functor $H : \calC^{op} \times \calC \to \calE$,
defined as an object in $\calE$ with arrows $\alpha_c : H(c,c) \to \int^{\calC}H$ for each $c \in \calC$
collectively universal with respect to
\[
\begin{tikzcd}
    H(c',c) \arrow[r, "f_*"] \arrow[d,"f^*"'] & H(c',c') \arrow[d, "\alpha_{c'}"] \\
    H(c,c) \arrow[r, "\alpha_c"] & \int^{\calC} H
\end{tikzcd}
\]

The coend is isomorphic to $$\colim(K/d \stackrel{U}{\to} \calC \stackrel{F}{\to} \calE)$$
where $K/d$ is the comma category with objects $(c, Kc \to d)$ and obvious morphisms
and $U$ is the forgetful functor.

Sometimes the category we care about, such that homotopy category, might not have all (co)limits,
(To be finished)

\begin{definition}
    A left Kan extension $(\Lan_K F, \eta)$ is \emph{preserved} by $L : \calE \to \calF$ if the
    whiskered composite $(L\Lan_K F, L \eta)$ is a left Kan extension of $LF$ along $K$.
\end{definition}

\begin{definition}(Def 1.3.4)
    A right Kan extension if \emph{pointwise} if it is preserved by representable functors $\calE(e,-)$ for all $e \in \calE$.
\end{definition}

\begin{remark}
    Pointwise left Kan extension is not simply defined as ``preserved by corepresentable'' since the variance does not match.
    Instead, a left Kan extension is pointwise if its corresponding right Kan extension
    under $(-)^{op}$ is a pointwise right Kan extension.
\end{remark}

\section{Homological algebra}\label{Homological-algebra}
(The following discussion applies to any \href{https://ncatlab.org/nlab/show/abelian+category}{abelian category},
but we will only focus on the origin of this concept $\Ab$, the catgeory of abelian groups.)

On the way to define \href{https://ncatlab.org/nlab/show/singular+homology}{singular homology} of a topological space $X$,
we obtain an intermediate product, the singular chain complexes $C_\bullet(X)$:
\[
    C_0(X) \stackrel{d_1}{\leftarrow} C_1(X) \stackrel{d_2}{\leftarrow} C_2(X)
    \cdots \stackrel{d_n}{\leftarrow} C_n(X) \stackrel{d_{n+1}}{\leftarrow} \cdots
\]
with the property that $d_i\circ d_{i+1} = 0$ for all $i$.
It follows that $\im(d_{i+1})$ is a subgroup of $\ker(d_i)$,
thus we can define the homology groups as
$$H_i(X) := \ker(d_{i}) / \im(d_{i+1}). $$

Forget about topology and abstract the notion, we can just say a (bounded below) chain complex of abelian groups is
a sequence of abelian groups $C_\bullet$:
\[
    C_0 \stackrel{d_1}{\leftarrow} C_1 \stackrel{d_2}{\leftarrow} C_2
    \cdots \stackrel{d_n}{\leftarrow} C_n \stackrel{d_{n+1}}{\leftarrow} \cdots
\]
such that $d_i \circ d_{i+1} = 0$ for all $i$.

A morphism between two chain complexes is a sequence of map between the component such that all the squares commute.
We denote the category of chain complexes of abelian group as $\ChAb$.

With homology, we can define the notion of quasi-isomorphisms: a map between $f_\bullet: C_\bullet \to D_\bullet$
such that $H_i(f) : H_i(C_\bullet) \to H_i(D_\bullet)$ is an isomorphism for all $i$.
Quasi-isomorphisms satisfy the 2-of-6 property, so they forms a class of weak equivalence.

Let $F$ be an endofunctor of $\Ab$, then it induces an endofunctor $F$ of $\ChAb$.
The bad thing is that it rarely preserves homology or preserves quasi-isomorphism.

Apply $-\otimes \bZ/2$ (see \href{https://ncatlab.org/nlab/show/tensor+product+of+abelian+groups}{this} for more) to
\[
    \cdots \longrightarrow 0 \longrightarrow \bZ \stackrel{\times 2}{\longrightarrow} \bZ \longrightarrow 0
\]
and we obtain
\[
    \cdots \longrightarrow 0 \longrightarrow \bZ/2 \stackrel{\times 2}{\longrightarrow} \bZ/2 \longrightarrow 0
\]
But the $H_2$ of the first chain complex is $0$ while the for the second it's $\bZ/2$.

$\Hom_{\bZ}(\bZ/2, -)$ (the right adjoint of $-\otimes \bZ/2$) doesn't preserve homology either (see Ex 2.1.15).

For the construction of the derived functor of $-\otimes A$, see Week \ref{part-3}.

\newpage
\part{Derived functors via deformations, featuring classical homological algebra --- Mar 17, led by}
\label{part-3}

\end{document}